\documentclass[12pt]{article}
\usepackage[utf8]{inputenc}
\usepackage[T1]{fontenc}
\usepackage[margin=1in]{geometry}
\usepackage{setspace}
\usepackage{xcolor}
\usepackage[colorlinks=true, urlcolor=blue!60!black, linkcolor=blue!60!black]{hyperref}
\setlength{\parindent}{2em}
\setlength{\parskip}{1.25em}
\renewcommand{\baselinestretch}{1.0}
\begin{document}
\begin{center}
{\LARGE \textit{Chiles v. Salazar}, 607 U.S. \_\_\_ (2026)\par}
\vspace{0.8em}
{\large Case Brief\par}
\end{center}
\vspace{1em}

\section{Memory Jogger}

In \href{https://www.courtlistener.com/opinion/10830655/chiles-v-salazar-revisions-33126/}{\textit{Chiles v. Salazar}}, No. 24-539, 607 U.S. \_\_\_ (2026), the Supreme Court held 8--1 that Colorado's Minor Conversion Therapy Law, as applied to a licensed counselor's pure talk therapy, discriminates based on viewpoint in violation of ordinary First Amendment rules and therefore triggers strict scrutiny rather than rational-basis review as a regulation of professional conduct.

\section{Detailed Case Facts}

Kaley Chiles holds a master's degree in clinical mental health and a Colorado state counseling license. In her verified complaint, Chiles alleged that she counsels clients on trauma, addiction, eating disorders, gender dysphoria, and sexuality without imposing predetermined goals, instead working with clients to pursue their own stated objectives. Some minor clients who are content with their sexual orientation or gender identity seek her help with family or social issues; others come to her hoping to ``reduce or eliminate unwanted sexual attractions, change sexual behaviors, or grow in the experience of harmony with [their] bod[ies].'' Chiles does not prescribe medication, use medical devices, or employ aversive or physical techniques---she provides exclusively talk therapy.

In 2019, Colorado joined roughly half the States in enacting a ban on conversion therapy for minors, titled the Minor Conversion Therapy Law (MCTL), Colo. Rev. Stat. \S~12-245-224(1)(t)(V) (2025). Colorado enacted the statute in response to concerns over youth mental health and medical consensus---including reports from the American Psychological Association (APA) and the Substance Abuse and Mental Health Services Administration (SAMHSA)---finding conversion therapy ineffective and associated with heightened depression, anxiety, and suicidality. Under Colo. Rev. Stat. \S~12-245-202(3.5)(a), ``conversion therapy'' is defined broadly to prohibit ``any practice or treatment \ldots\ that attempts \ldots\ to change an individual's sexual orientation or gender identity,'' including ``efforts to change behaviors or gender expressions or to eliminate or reduce sexual or romantic attraction or feelings toward individuals of the same sex.'' At the same time, \S~12-245-202(3.5)(b)(I)--(II) explicitly permits practices providing ``[a]cceptance, support, and understanding for the facilitation of an individual's \ldots\ identity exploration and development'' as well as ``[a]ssistance to a person undergoing gender transition.'' Any Coloradan may file a regulatory complaint (\S~12-245-226(1)(a)(II)(B)), exposing a licensed counselor to fines, probation, or license revocation (\S~12-245-225).

Chiles brought a pre-enforcement, as-applied First Amendment challenge in federal court against Patty Salazar, Executive Director of the Colorado Department of Regulatory Agencies, and state licensing officials. Chiles did not contest Colorado's authority to ban aversive physical interventions (such as electric shock therapy); she challenged the MCTL solely as applied to her voluntary, speech-only talk therapy. The case reached the Supreme Court amid a deep circuit split over whether talk-therapy conversion bans regulate professional conduct subject to rational-basis review (as the Ninth and Tenth Circuits had held) or regulate speech subject to heightened First Amendment scrutiny (as the Eleventh and Sixth Circuits had held), following the Court's landmark rejection of a diminished ``professional speech'' doctrine in \textit{National Institute of Family and Life Advocates v. Becerra} (\textit{NIFLA}), 585 U.S. 755 (2018) (\href{https://constitution.congress.gov/browse/essay/intro.9-4-1/ALDE_00000126/}{Constitution Annotated}).

\section{Procedural History}

Chiles filed suit in the United States District Court for the District of Colorado seeking a preliminary injunction against enforcement of the MCTL as applied to her talk therapy. On December 19, 2022, the District Court (Civ. Action No. 1:22-cv-2287) held that Chiles had Article III pre-enforcement standing because she faced a credible threat of enforcement, but denied a preliminary injunction on the merits, concluding that the MCTL regulated professional conduct and burdened speech only incidentally, thus triggering rational-basis review.

The United States Court of Appeals for the Tenth Circuit affirmed in a 2--1 decision, 116 F.4th 1178 (10th Cir. 2024). The panel majority agreed that Chiles had Article III standing, but held that the MCTL permissibly regulated ``treatments'' and ``therapeutic modalities'' as professional conduct, satisfying rational-basis review. Judge Hartz dissented, arguing that recasting pure talk therapy as ``professional conduct'' was ``a labeling game'' barred by Supreme Court precedent.

Chiles petitioned for certiorari to resolve the circuit split (\href{https://www.scotusblog.com/cases/chiles-v-salazar/}{SCOTUSblog Case Page}). The Supreme Court granted certiorari, 604 U.S. 1193, heard oral argument on October 7, 2025 (\href{https://www.oyez.org/cases/2025/24-539}{Oyez}), and on March 31, 2026, reversed and remanded in an 8--1 decision (\href{https://www.supremecourt.gov/opinions/25pdf/24-539_fd9g.pdf}{Official Slip Opinion}).

\section{Judicial Votes}

\textbf{Majority (8):} Justice Gorsuch authored the opinion of the Court, joined in full by Chief Justice Roberts and Justices Thomas, Alito, Sotomayor, Kagan, Kavanaugh, and Barrett.

\textbf{Concurrence (2):} Justice Kagan filed a concurring opinion, joined by Justice Sotomayor.

\textbf{Dissent (1):} Justice Jackson filed a solo dissenting opinion.

\section{Holding}

\textbf{Judgment:} Reversed and remanded (116 F.4th 1178 reversed and remanded for further proceedings under the proper, heightened First Amendment standard of scrutiny).

\textbf{Rule Announced:} A state licensing law prohibiting licensed mental-health counselors from engaging in conversion therapy, as applied to pure talk therapy, directly regulates speech based on content and viewpoint---not conduct with an incidental burden on speech---and therefore triggers strict First Amendment scrutiny rather than rational-basis or intermediate review. States cannot evade strict scrutiny by relabeling pure speech as a ``treatment'' or ``therapeutic modality,'' nor by aggregating licensing, informed-consent, and malpractice traditions into a new ``First Amendment Free Zone'' for speech deemed ``substandard care.''

\textbf{Holding vs. Dicta:} The Court's binding holdings are (1) Chiles possessed Article III pre-enforcement standing (confirmed in the majority's Footnote~*), and (2) Colorado's MCTL, as applied to Chiles's speech-only talk therapy, is a viewpoint-based restriction on speech subject to strict scrutiny, requiring reversal of the Tenth Circuit's rational-basis ruling. Crucially, the Court did \textit{not} itself apply strict scrutiny or definitively strike down the MCTL on the merits, instead remanding for the lower courts to conduct that inquiry. Nor did the Court resolve the constitutionality of aversive physical conversion therapies (which Chiles conceded were valid) or of content-based, viewpoint-neutral regulations of healthcare speech (explicitly flagged as an open question by Justice Kagan's concurrence).

\section{Analysis of Opinions}

\subsection{Majority Opinion (Justice Gorsuch)}

Writing for an eight-Justice majority, Justice Gorsuch framed the case around bedrock First Amendment hostility to viewpoint discrimination (\href{https://www.supremecourt.gov/opinions/25pdf/24-539_fd9g.pdf}{Slip Op. at 1--23}). First, in Footnote~*, the Court affirmed Chiles's Article III standing under \textit{Susan B. Anthony List v. Driehaus}, 573 U.S. 149 (2014), rejecting Colorado's late-breaking, narrowed reading of the MCTL as inconsistent with the statutory text and noting that Colorado had litigated the case for three years while expressly declining at oral argument to disavow enforcement against Chiles.

On the merits, Gorsuch began with first principles: content-based speech restrictions are presumptively unconstitutional and trigger strict scrutiny (\textit{Reed v. Town of Gilbert}, 576 U.S. 155, 163 (2015)), and viewpoint discrimination is an even more ``egregious form'' of content regulation (\textit{Rosenberger v. Rector and Visitors of Univ. of Va.}, 515 U.S. 819, 829 (1995)). Because Chiles provides exclusively talk therapy, ``all Ms. Chiles does is speak---and, as far as she is concerned, speech is all Colorado seeks to regulate.'' Under the MCTL, she may express ``[a]cceptance, support, and understanding'' for identity exploration or assist a minor undergoing gender transition, yet she is forbidden from speaking in ways that help a consenting minor seek to reduce unwanted attractions or realign gender identity with biological sex. Relabeling spoken words as ``conduct,'' ``treatment,'' or a ``therapeutic modality'' cannot evade the First Amendment (\textit{NAACP v. Button}, 371 U.S. 415, 429 (1963)).

Gorsuch then rejected Colorado's reliance on the speech-incident-to-conduct exception. Under \textit{Cohen v. California}, 403 U.S. 15 (1971), and \textit{Holder v. Humanitarian Law Project}, 561 U.S. 1 (2010), a law that generally regulates conduct (such as disturbing the peace or banning aversive physical therapies) still triggers strict scrutiny when applied to pure speech. Properly understood, the speech-incident-to-conduct doctrine applies only when speech is integrally bound up with separately unlawful conduct (like selling contraband under \textit{United States v. Williams}, 553 U.S. 285 (2008)) or when a content-neutral conduct ban incidentally sweeps in expression (like an outdoor fire ban applied to flag burning). Colorado's law satisfied neither test: non-licensees such as life coaches could lawfully give the exact same advice, and the MCTL trained directly on Chiles's viewpoint.

Finally, Gorsuch rejected Colorado's bid under \textit{NIFLA} and \textit{United States v. Stevens}, 559 U.S. 460 (2010), to recognize a new historic tradition exempting ``substandard care'' from strict scrutiny. First, aggregating licensing, informed consent, and malpractice operated at too high a level of generality. Second, taken individually, each tradition failed: (1) state counselor-licensure statutes date only to 1976 and licensing traditionally governs qualifications rather than viewpoints; (2) informed-consent laws (\textit{Planned Parenthood of Southeastern Pa. v. Casey}, 505 U.S. 833 (1992)) attach to separate physical medical procedures and require factual, uncontroversial disclosures; and (3) medical malpractice requires exacting proof of causation and injury---providing ``breathing room for protected speech''---and sometimes permits patient consent to nonstandard care, whereas the MCTL imposes administrative penalties simply for expressing a prohibited view. Warning that medical consensus is not static---citing the APA's 1952 classification of homosexuality as a mental disorder and \textit{Buck v. Bell}, 274 U.S. 200 (1927)---Gorsuch concluded that the First Amendment forbids turning prevailing professional opinion into enforced orthodoxy.

\subsection{Concurring Opinion (Justice Kagan, joined by Justice Sotomayor)}

Justice Kagan joined the Court's opinion in full but wrote separately to underscore the decision's narrow reach (\href{https://www.supremecourt.gov/opinions/25pdf/24-539_fd9g.pdf}{Kagan, J., concurring, at 1--4}). Kagan agreed that Colorado's MCTL was ``textbook'' viewpoint discrimination: it allowed therapists to affirm a minor's gay or transgender identity while forbidding them from helping a minor try to change it. Pointing to the mirror-image hypothetical---a state law banning affirming talk therapy while permitting conversion therapy---Kagan noted that the First Amendment must apply identically regardless of which side the State favors.

Kagan wrote primarily to reserve a different question: how courts should evaluate \textit{content-based but viewpoint-neutral} regulations of healthcare professionals' speech. Citing \textit{Davenport v. Washington Education Assn.}, 551 U.S. 177 (2007), \textit{Vidal v. Elster}, 602 U.S. 286 (2024), and her own concurrence in \textit{Reed}, Kagan stressed that strict scrutiny is relaxed when a content-based law poses no realistic danger of ``official suppression of ideas.'' Because medical regulation inevitably touches speech, the distinction between viewpoint-based and viewpoint-neutral content rules could prove ``decisive'' in future healthcare cases. In a footnote, Kagan pushed back against Justice Jackson's claim that no such category exists, pointing out that many of the general competence and client-welfare regulations Jackson feared were endangered are content-based yet viewpoint-neutral.

\subsection{Dissenting Opinion (Justice Jackson)}

In a 35-page solo dissent, Justice Jackson argued that the majority ignored the ``crucial context'' of state medical regulation and misread \textit{NIFLA} and \textit{Casey} (\href{https://www.supremecourt.gov/opinions/25pdf/24-539_fd9g.pdf}{Jackson, J., dissenting, at 1--35}). Opening with \textit{Lambert v. Yellowley}, 272 U.S. 581, 596 (1926) (``[T]here is no right to practice medicine which is not subordinate to the police power of the States''), Jackson detailed the overwhelming medical consensus from the APA and SAMHSA that conversion therapy---including nonaversive talk therapy---is ineffective and inflicts severe psychological harm, shame, and suicidality on minors.

Doctrinally, Jackson contended that \textit{NIFLA}'s discussion of \textit{Casey} turned not on whether a physical scalpel or drug was involved, but on the State's \textit{objective}: whether the State is regulating ``speech as speech'' to suppress an idea, or ``incidentally'' restricting speech in the course of regulating a medical treatment tied to a patient-provider relationship. Unlike the waiting-room notice in \textit{NIFLA} (which was not tied to any medical procedure), the MCTL governs only clinical treatment within a licensed provider-patient relationship, leaving Chiles completely free to advocate for conversion therapy in speeches, articles, or public debate. Drawing on scholarship by Claudia Haupt and Robert Post, Jackson argued that ``professional medical speech'' does not serve the marketplace of ideas---in clinical care, ``we insist upon competence, not debate''---and that medical standards of care are \textit{necessarily} viewpoint-based because States must choose safe, effective treatments over harmful ones (such as prohibiting therapists from encouraging suicide or telling anorexic patients to eat less).

Finally, Jackson reviewed centuries of medical licensing (dating to 1639) and common-law malpractice liability (including liability for harmful medical advice) to show that States have always regulated speech-based medical care. She warned that the majority's ruling ``plays with fire'' and threatens to render speech-only therapies unregulatable, potentially immunizing cruel, incompetent, or harmful counseling practices from state licensing discipline.

\section{Examples: Future Applications}

\begin{enumerate}
  \item \textbf{Same-Side Application (State Ban on Gender-Affirming Talk Therapy for Minors):} A State enacts a licensing statute prohibiting licensed therapists from providing talk therapy that affirms a minor's transgender identity or supports social transition, while expressly permitting counseling aimed at aligning the minor's gender identity with biological sex. \textit{Outcome under Chiles:} Subject to strict scrutiny and presumptively unconstitutional. Both Justice Gorsuch's majority opinion and Justice Kagan's concurrence explicitly confirm that the mirror-image law discriminates based on viewpoint in the exact same way as Colorado's MCTL.

  \item \textbf{Same-Side Application (State Prohibition on Counselors Discussing Abortion or Assisted-Suicide Options):} A State prohibits licensed mental-health counselors from speaking favorably about abortion or medical aid in dying during talk-therapy sessions, while allowing counselors to encourage carrying pregnancies to term or declining aid in dying. \textit{Outcome under Chiles:} Triggers strict scrutiny. Because the counseling consists solely of spoken dialogue unattached to a physical medical procedure performed by the counselor, and the statute favors one viewpoint over another, the State cannot shield it as ``professional conduct.''

  \item \textbf{Opposite-Side Application (Ban on Aversive Physical Conversion Techniques or Unapproved Pharmaceuticals):} A State prohibits licensed practitioners from using electric shocks, nausea-inducing drugs, or unapproved hormone regimens to attempt to alter a minor's sexual orientation or gender identity. \textit{Outcome under Chiles:} Valid under rational-basis review (no First Amendment strict scrutiny). As Gorsuch repeatedly noted, Chiles conceded that the MCTL has valid applications to physical interventions and medications, which constitute non-expressive conduct squarely within the State's traditional police power.

  \item \textbf{Opposite-Side Application (Factual Informed-Consent Disclosures Prior to Surgery):} A State requires surgeons to disclose objective, noncontroversial medical risks and alternatives to a patient before performing an elective surgical procedure. \textit{Outcome under Chiles:} Exempt from strict scrutiny. Both the majority and \textit{NIFLA} reaffirm that informed-consent requirements tied to a separate physical procedure that would otherwise constitute battery---and limited to factual, uncontroversial information---fall within the recognized speech-incident-to-conduct and commercial-disclosure exceptions.

  \item \textbf{Fence-Sitter Application (Viewpoint-Neutral Ban on ``Predetermined-Outcome'' Therapy or General ``Incompetence'' Discipline):} Following \textit{Chiles}, a State amends its statute (as Colorado itself moved to do in May 2026) to prohibit licensed therapists from ``directing a patient toward \textit{any} predetermined sexual orientation or gender identity outcome''---whether heterosexual, gay, cisgender, or transgender---or disciplines a talk therapist under a general rule against ``incompetent'' or ``cruel'' care. \textit{Analysis:} This is the exact fault line between the concurrence and dissent. Under Justice Kagan's concurrence (joined by Justice Sotomayor), a content-based yet viewpoint-neutral rule that does not favor one ideological side over the other might escape strict scrutiny under \textit{Vidal v. Elster} and \textit{Davenport}. Under Justice Jackson's dissent, however, enforcing a ``best interests'' or ``incompetence'' rule against a talk therapist inevitably judges the clinician's viewpoint in application, and conservative Justices committed to \textit{Reed}'s strict-scrutiny baseline for content-based rules may still strike down such statutes unless tightly tethered to proof of actual patient injury or fraud.
\end{enumerate}

\section{Critique}

\textit{Chiles v. Salazar} immediately generated intense scholarly debate, both for its lopsided 8--1 alignment and for its sweeping implications for healthcare regulation (\href{https://www.scotusblog.com/2026/04/the-non-partisan-puzzle-in-the-conversion-therapy-case/}{Craig Konnoth, SCOTUSblog}; \href{https://www.naag.org/attorney-general-journal/supreme-court-report-volume-33-issue-10/}{NAAG Supreme Court Report}).

\textbf{Progressive and Health-Law Perspective:} Progressive and health-law scholars---echoing Justice Jackson's dissent and Professor Craig Konnoth's SCOTUSblog critique---argue that the majority's ``speech versus physical procedure'' boundary is clinically incoherent. Enormous swaths of modern medicine consist entirely of spoken diagnosis and advice untethered to a physical procedure (e.g., advising ``watchful waiting'' for pediatric infections, counseling dementia families, or delivering psychotherapy). Critics also challenge Gorsuch's distinction between administrative licensing bans and medical malpractice tort suits: in Colorado, both civil malpractice actions and administrative licensing proceedings generally apply a preponderance-of-the-evidence standard, calling into question the majority's claim that malpractice uniquely supplies ``exacting proof requirements.'' Progressive commentators further note the irony of the majority invoking \textit{Buck v. Bell}---whose eugenic premises historically targeted gay people as ``defective''---to protect a practice that every major U.S. medical association condemns as harmful to LGBTQ+ youth. At the same time, some observers interpret Justices Kagan and Sotomayor's decision to join the majority as a strategic compromise designed to lock in symmetrical First Amendment protection for gender-affirming talk therapy and to preserve room for viewpoint-neutral clinical regulations.

\textbf{Originalist, Textualist, and Free-Speech Libertarian Perspective:} Free-speech advocates and originalists largely praise \textit{Chiles} as a principled vindication of \textit{NIFLA}, \textit{Reed}, and \textit{Rosenberger}. On this view, allowing States to invoke ``medical consensus'' to silence disfavored viewpoints in voluntary talk therapy would give legislatures a blank check to enforce ideological orthodoxy on contested social and ethical questions---precisely the danger illustrated by past medical consensuses that pathologized homosexuality or endorsed eugenics. Yet some strict originalists still see unresolved tensions in \textit{Chiles}: the majority's reliance on modern strict-scrutiny tiering and its narrow framing of the historical inquiry (focusing on 1976 counselor-licensure laws rather than broader Founding-era police-power regulation of healers) leaves unclear how lower courts should reconcile the First Amendment with traditional ex ante licensing rules against fraudulent, abusive, or demonstrably dangerous verbal therapies without either gutting patient protection or diluting strict scrutiny.

\section{Key Quotations}

\begin{quote}
``Her speech does not become conduct just because the State may call it that. Nor does her speech become conduct just because it can also be described as a `treatment,' a `therapeutic modality,' or anything else. The First Amendment is no word game. And the rights it protects cannot be renamed away or their protections nullified by `mere labels.' '' --- Justice Gorsuch, Opinion of the Court, slip op. at 12--13.
\end{quote}

\begin{quote}
``Medical consensus, too, is not static; it evolves and always has. A prevailing standard of care may reflect what most practitioners believe today, but it cannot mark the outer boundary of what they may say tomorrow. Far from a test of professional consensus, the First Amendment rests instead on a simple truth: `[T]he people lose' whenever the government transforms prevailing opinion into enforced conformity.'' --- Justice Gorsuch, Opinion of the Court, slip op. at 22.
\end{quote}

\begin{quote}
``But the First Amendment stands as a shield against any effort to enforce orthodoxy in thought or speech in this country. It reflects instead a judgment that every American possesses an inalienable right to think and speak freely, and a faith in the free marketplace of ideas as the best means for discovering truth. However well-intentioned, any law that suppresses speech based on viewpoint represents an `egregious' assault on both of those commitments.'' --- Justice Gorsuch, Opinion of the Court, slip op. at 23.
\end{quote}

\begin{quote}
``And that means the `difference between viewpoint-based and viewpoint-neutral content discrimination' in the health-care context could prove `decisive.' Fuller consideration of that question, though, can wait for another day. We need not here decide how to assess viewpoint-neutral laws regulating health providers' expression because, as the Court holds, Colorado's is not one.'' --- Justice Kagan, Concurring, slip op. at 4.
\end{quote}

\begin{quote}
``The Constitution does not pose a barrier to reasonable regulation of harmful medical treatments just because substandard care comes via speech instead of scalpel.'' --- Justice Jackson, Dissenting, slip op. at 35.
\end{quote}

\end{document}
